\pagestyle{plain} % No headers, just page numbers

\chapter{\uppercase{Conclusion}} \label{chap:conclusion}

This dissertation addresses a fundamental scalability bottleneck in combinatorial
optimization on graphs: the mismatch between enormous ground sets and small optimal
solutions. We proposed a coherent framework of pruning strategies that reduce the ground
set to a much smaller subset before invoking a downstream solver, enabling practical
scalability without sacrificing solution quality. Our contributions span three
complementary directions: provably correct pruning, learned pruning without handcrafted
features, and automated feature discovery with constraint adaptability.

\textbf{QuickPrune}~\citep{pmlr-v258-nath25a} establishes the first theoretical
foundation for search space pruning in combinatorial optimization. Prior to this work,
all pruning methods were heuristics with no formal performance guarantees. QuickPrune is
a lightweight, single-pass streaming algorithm that processes each element exactly once
and produces a pruned ground set of size logarithmic in $n$, while retaining a provable
fraction of the optimal value. The algorithm operates under $\gamma$-weak submodularity,
a condition satisfied by a broad class of objectives including Maximum Coverage and
Influence Maximization, and handles both size and knapsack constraints. Empirically,
QuickPrune prunes over 90\% of the ground set while outperforming or matching all
classical and learning-based baselines across three CO problems and eight large-scale
datasets.

\textbf{H-DeepPruner}~\citep{nath2025hdeeppruner} demonstrates that effective learned
pruning does not require handcrafted, problem-specific features. All prior learning-based
pruners depend on domain knowledge to construct input features, limiting their
generality. H-DeepPruner replaces engineered features with random node initializations
and uses a two-stage pipeline: a GNN pre-pruner coarsely filters nodes in a single pass,
followed by a Neural MCTS fine-pruner that captures element interdependencies by
building the pruned set from scratch. Despite using no domain knowledge, H-DeepPruner
matches or outperforms carefully engineered baselines on real-world graphs, establishing
that graph topology alone suffices for competitive learned pruning.

\textbf{LLM2Prune} closes the remaining gap by automating feature discovery and
enabling natural adaptation to diverse constraints. Rather than relying on random or
handcrafted features, LLM2Prune uses large language models to propose node-level
features from a task description, guided by an LLM-directed beam search over the feature
space. Performance evaluation and feature-importance scores from GNNExplainer serve as
feedback signals, steering the search toward high-quality, problem-grounded features.
Switching from size to knapsack constraints requires only changing the prompt, with no
algorithmic redesign. Across three CO problems and eight datasets, LLM2Prune is
competitive with classical methods and outperforms all learning-based baselines, while
completing inference in under three seconds even on graphs with over one million nodes.

Together, these three contributions form a coherent and progressively more capable
framework for scalable combinatorial optimization on graphs. QuickPrune provides
theoretical guarantees; H-DeepPruner shows that learning-based pruning can succeed
without domain knowledge; and LLM2Prune automates the feature discovery process and
generalizes across constraint types.
